\section{Experimental Setup}
\label{sec:setup}

This section documents the single-GPU pipeline used to produce all subsequent results. The pipeline is released in full at \url{https://github.com/Euswbnix/Machine_translation} (pretraining) and \url{https://github.com/Euswbnix/Machine-Translation-SFT} (post-training).

\subsection{Hardware and software}
All training runs use a single NVIDIA RTX 5090 (Blackwell, sm\_120, 32~GB VRAM, 600~W TDP). Software stack: Python 3.10, PyTorch 2.8 nightly (CUDA 12.8 build for sm\_120), SentencePiece 0.1.99, sacrebleu 2.3, and CometKiwi-22 via \texttt{unbabel-comet 2.2}. We train in bf16 throughout (autocast); a brief fp16 attempt on Big diverged and was abandoned.

\subsection{Data}
\paragraph{WMT14 en-fr.} The raw parallel corpus contains $\sim$~40.8M sentence pairs (Europarl, News Commentary, Common Crawl, UN, Giga). We construct two cleaned variants from this same source:
\begin{itemize}
\item \textbf{v1 (strict, 9.3M).} The raw corpus is capped at 10M pairs (the first 10M from the HuggingFace Datasets stream, deterministic order), then cleaned with: empty/duplicate removal, length range [3, 200] tokens, src/tgt length ratio $\in [0.5, 2.0]$, and Latin-script ratio $\geq 0.9$. 9.3M pairs survive (93\% keep rate on the 10M head).
\item \textbf{v2 (loose, 30M).} The full 40.8M raw corpus, cleaned with identical per-pair filters but without the 10M cap. 30.1M pairs survive (74\% keep rate over all 40.8M; lower because the long tail beyond 10M contains more noisy or misaligned pairs).
\end{itemize}
The v1 / v2 split forms the data-quality axis of the 2x2 ablation: identical per-pair filters, different breadth.

\paragraph{WMT14 en-de.} Raw $\sim$~4.5M pairs (Europarl, News Commentary, Common Crawl). The same cleaning thresholds as v1 are applied (default settings of our cleaning script). 4.17M pairs survive (92.6\% keep rate; the en-de raw corpus is cleaner per-pair than en-fr).

\paragraph{Eval splits.} newstest2013 (3,000 pairs, valid) and newstest2014 (3,003 pairs, test) for both en-fr and en-de, downloaded via the WMT14 HuggingFace dataset.

\subsection{Tokenization}
SentencePiece BPE~\cite{sennrich2016neural, kudo2018sentencepiece}, joint en-fr (or en-de) vocabulary of 32{,}000 subwords, with \texttt{character\_coverage=1.0} so that accented characters (\textit{Israël}, \textit{Noël}, German umlauts) are not mapped to \texttt{<unk>}. The SPM is trained on the cleaned training corpus; valid and test sides reuse the same SPM. We share the input embedding, output embedding, and pre-softmax projection (Vaswani's \emph{tied embeddings} option).

\subsection{Model architectures}
We use the original Vaswani configurations. Base ($d_\text{model}{=}512$, $h{=}8$, $d_\text{ff}{=}2048$, $L{=}6{+}6$) totals 60.5M parameters with shared embeddings; Big ($d_\text{model}{=}1024$, $h{=}16$, $d_\text{ff}{=}4096$, $L{=}6{+}6$) totals 209M. Dropout is 0.1 throughout (en-fr), or 0.3 for en-de Big (following the original paper's stronger regularization for the smaller corpus). Maximum sequence length is 256 tokens.

\subsection{Training}
\paragraph{Optimization.} AdamW ($\beta_1{=}0.9$, $\beta_2{=}0.98$, $\epsilon{=}1\text{e-}9$), label-smoothed cross-entropy ($\alpha{=}0.1$), gradient clipping at $\|g\|{=}1.0$. Noam learning-rate schedule~\cite{vaswani2017attention} with $d_\text{model}^{-0.5}$ scale and a min-LR floor of $1\text{e-}5$ to prevent the late-Noam tail from approaching zero. Warmup 4K steps for Base / 8K for Big.

\paragraph{Effective batch.} We use token-based dynamic batching: each micro-batch packs $\leq$~24{,}576 tokens (Base) / 8{,}192 (Big), with $\leq$~192 / 64 sentences per micro-batch, and we accumulate 4 micro-batches per optimizer step. The effective batch is therefore $\sim$~98K tokens/step (Base) / 32K tokens/step (Big), matching the original paper's $\sim$~25K tok/batch on en-de once normalized by device count.

\paragraph{Run lengths.} Base v1.1 trains for 122K steps ($\sim$~3 epochs), Big v1.1 for 260K (early-stopped from a 280K plan), Base v2 for 600K, Big v2 for 416K (halted), Base en-de for 230K. Each Base run is $\sim$~2--3~h on the 5090; each Big run is $\sim$~6--7~h.

\paragraph{Loss-spike guard.} Each accumulated batch's average loss is compared against an EMA of recent batches (window $\alpha{=}0.005$, $\sim$~200 steps). If the current batch exceeds $1.3\times$ the EMA, the entire effective batch's gradient is discarded. The intent is to drop poisoned micro-batches without forcing a hard halt. Trigger counts are saved to the training report (\S\ref{sec:methodology}).

\paragraph{Adaptive eval interval.} Validation runs every 20K steps when the smoothed training loss is high (above 4.0) and every 2--3K steps once it crosses below 2.8, with linear interpolation in between. This avoids wasting compute on early-training BLEU evals (which are essentially noise) while keeping the late-training signal fine-grained.

\subsection{Evaluation}
\paragraph{BLEU.} sacrebleu 13a (default \texttt{intl} normalization, default \texttt{13a} tokenizer)~\cite{post2018call}. We report newstest2013 (valid) and newstest2014 (test). For en-de we additionally try the original paper's beam=4, length-penalty=0.6 setting (\S\ref{sec:ende}).

\paragraph{Decoding.} Beam search, beam size 5, length penalty 1.0 (en-fr) / 0.6 (en-de when comparing to the original paper), maximum decoded length 256.

\paragraph{Checkpoint averaging.} The classic Vaswani trick: average the parameters of the last 5 checkpoints (5K-step intervals near convergence) before evaluation. This consistently lifts BLEU by 0.3--0.5. All headline numbers in the paper are post-averaging unless noted otherwise.

\paragraph{Single-checkpoint comparison.} For fine-tuning runs (\S\ref{sec:sft}), the run is too short to produce 5 save-interval-spaced checkpoints, so we evaluate the final single checkpoint. We flag this unit mismatch in Table~\ref{tab:sft_results} and note that it inflates the apparent fine-tuning degradation by roughly the averaging gain (0.3--0.5 BLEU).

\subsection{Compute budget}
The full study --- five headline pretraining runs (en-fr Base v1.1, Big v1.1, Base v2, Big v2; en-de Base), two fine-tuning runs, and 13.8~h of CometKiwi scoring --- totals $\sim$~30 GPU-hours on the 5090.
